\documentclass[12pt,oneside]{article}

\usepackage[margin=1in]{geometry}
\usepackage{setspace}
\usepackage{fancyhdr}
\usepackage{sectsty}
\usepackage{enumitem}
\usepackage[hidelinks]{hyperref}
\usepackage{times}

\onehalfspacing

\setlist[itemize]{leftmargin=0.5in}
\setlist[enumerate]{leftmargin=0.5in}

\pagestyle{fancy}
\fancyhf{}
\fancyhead[C]{CRIMINAL INVESTIGATION REFERRAL: INSTITUTIONAL CORRUPTION \& CIVIL RIGHTS VIOLATIONS}
\fancyfoot[C]{\thepage}
\renewcommand{\headrulewidth}{0.5pt}

\sectionfont{\normalsize\bfseries}
\subsectionfont{\normalsize\bfseries\underline}

\begin{document}

\begin{center}

\vspace*{0.5in}

{\textbf{\Large COMPREHENSIVE CRIMINAL INVESTIGATION REFERRAL}}

\vspace{0.2in}


{\textbf{DUAL PERPETRATOR CONSPIRACY}}

\vspace{0.1in}


{\textbf{INSTITUTIONAL CORRUPTION \& ENDEMIC CIVIL RIGHTS VIOLATIONS}}

\vspace{0.1in}


{\textbf{Danielle Chard, E'Lise Chard, YWCA Missoula,}}

\vspace{0.05in}


{\textbf{Missoula Police Department, and Associated Officials}}

\vspace{0.1in}


2015--2025

\vspace{0.5in}

\end{center}

\begin{flushleft}

\textbf{TO:} \\
Federal Bureau of Investigation --- Civil Rights Division, Seattle Field Office \\
Federal Bureau of Investigation --- Civil Rights Division, Montana Field Office \\
U.S. Attorney, District of Montana \\
U.S. Attorney, Western District of Washington \\
U.S. Department of Justice --- Criminal Division, Public Integrity Section \\
Montana Attorney General --- Criminal Division \\
Washington Attorney General --- Civil Rights Division

\vspace{0.3in}

\textbf{FROM:} \\
MISJustice Alliance \\
On behalf of Elvis Ryland Nuno and other victims of institutional corruption

\vspace{0.3in}

\textbf{DATE:} December 8, 2025

\vspace{0.3in}

\textbf{RE:} Comprehensive Criminal Investigation Referral \\
Pattern of Racketeering, Civil Rights Violations, Institutional Corruption \\
Federal RICO, 18 U.S.C. § 1962; Civil Rights Conspiracy, 42 U.S.C. § 1985 \\
Interstate Criminal Activity (Washington \& Montana, 2015--2025)

\end{flushleft}

\newpage

\section*{CASE SUMMARY}

MISJustice Alliance respectfully submits this \textbf{comprehensive criminal investigation referral} documenting a decade-long pattern of institutional corruption, civil rights violations, and coordinated criminal activity spanning Washington State and Montana (2015--2025). While the majority of documentation focuses on the systematic victimization of \textbf{Elvis Ryland Nuno}, the evidence establishes that individual criminal conduct by \textbf{Danielle Christine Chard} and \textbf{E'Lise Marie Chard} represents merely the visible manifestation of endemic institutional corruption that is part of a larger and more substantial pattern of which \textbf{the YWCA of Missoula} sits at the nexus through which numerous individuals have been victimized, and continues unchecked.

This submission requests immediate federal and state criminal investigations targeting:

\begin{enumerate}
\item \textbf{Individual Perpetrators:} Danielle Chard and E'Lise Chard
\item \textbf{Institutional Enterprise:} YWCA Missoula and associated governance
\item \textbf{Law Enforcement Actors:} Missoula Police Department officers and supervisors
\item \textbf{Prosecutorial Actors:} Missoula County and King County prosecutors
\item \textbf{Pattern of Endemic Corruption:} Systematic abuse extending beyond reasonable doubt to warrant expanded investigation of institutional patterns affecting multiple victims over more than a decade
\end{enumerate}

\textbf{YWCA Notice and Non-Response:} MISJustice Alliance sent formal notices to YWCA Missoula on October 12, 2025 and December 5, 2025, documenting systemic issues, conflicts of interest, and the institution's role in the coordinated harassment campaign against Mr. Nuno.\footnote{See Evidentiary Documentation \S{} 2.9--2.10, \texttt{Formal-Complaints/ywca-first-notice-10-12-2025.pdf} and \texttt{ywca-second-notice-12-5-2025.pdf}.} These notices specifically identified Detective Connie Brueckner's dual role as law enforcement officer and YWCA board member, detailed patterns of retaliation and confidential information misuse, and requested explanations of institutional practices. \textit{To date, YWCA Missoula has provided no response to either notice}, demonstrating continued institutional indifference to serious allegations of constitutional violations and establishing a pattern of deliberate refusal to address complaints---the same pattern that enabled the decade-long harassment of Mr. Nuno following his 2018 formal complaint.\footnote{This non-response, following detailed legal notices copied to city and county officials, supports claims of deliberate indifference under \textit{Monell v. Department of Social Services}, 436 U.S. 658 (1978), and consciousness of institutional guilt.}

\section*{SUBMISSION CONTENTS}

This comprehensive referral package includes eleven (11) components:

\begin{enumerate}
\item \textbf{This Case Summary} --- Overview of criminal enterprise and referral scope
\item \textbf{FBI Seattle Civil Rights Cover Letter} --- Federal jurisdiction and RICO/civil rights violations \& interstate coordination originating from Washington
\item \textbf{FBI Montana Civil Rights Cover Letter} --- Federal jurisdiction and RICO/civil rights violations \& interstate coordination originating from Montana
\item \textbf{Montana Attorney General Cover Letter} --- State criminal prosecution referral
\item \textbf{Washington Attorney General Cover Letter} --- Interstate coordination and Washington-based misconduct
\item \textbf{U.S. Attorney, District of Montana Cover Letter} --- Federal RICO and civil rights prosecution (Montana)
\item \textbf{U.S. Attorney, Western District of Washington Cover Letter} --- Federal civil rights and corruption prosecution (Washington)
\item \textbf{DOJ Public Integrity Section Cover Letter} --- Public corruption and prosecutorial misconduct
\item \textbf{Danielle Chard Criminal Report} --- Individual perpetrator analysis and charges
\item \textbf{E'Lise Chard Criminal Report} --- Individual perpetrator analysis and institutional nexus
\item \textbf{YWCA Institutional Corruption Supplemental; Elvis Nuno Sworn Declaration; Comprehensive Evidentiary Documentation}
\end{enumerate}

\section*{CRITICAL FINDINGS WARRANTING IMMEDIATE INVESTIGATION}
\textbf{Research Methodology:} This criminal referral represents the collaborative product of over 1,000 hours of systematic research, legal analysis, and evidence organization by the MISJustice Alliance legal research collective. Given the multi-jurisdictional nature of this case (Washington and Montana, 2015--2025), extensive cross-referencing of court records, police reports, witness statements, and institutional communications was necessary to establish the coordinated pattern of civil rights violations. All evidence has been indexed, authenticated, and organized to facilitate efficient investigation and prosecution.


\subsection*{Individual Criminal Conduct (Chard Sisters)}

\textbf{Danielle Christine Chard} (DOB: 09/23/1983):
\begin{itemize}
\item Interstate false reporting pattern (2015--2025)
\item Criminal defamation across multiple jurisdictions
\item Perjury in sworn statements and court filings
\item Witness tampering and evidence fabrication
\item Coordinated harassment campaign with estranged sister E'Lise
\item Manipulation of law enforcement through false victimization claims
\end{itemize}

\textbf{E'Lise Marie Chard} (DOB: 09/23/1983):
\begin{itemize}
\item Institutional conspiracy through YWCA employment
\item Perjury in police reports and court proceedings
\item Criminal defamation and false light invasion of privacy
\item Coordination with law enforcement (Detective Connie Brueckner)
\item First Amendment retaliation following formal YWCA complaint
\item Professional position abuse to weaponize institutional resources
\item Intentional manipulation of subordinates and vulnerable clients
\end{itemize}

\subsection*{Institutional Corruption (YWCA Missoula)}

The YWCA Institutional Corruption Supplemental (Document \#8) establishes RICO predicate acts \textbf{independent of the Nuno case}, demonstrating systematic criminal enterprise activity:

\begin{itemize}
\item \textbf{Governance Capture:} Detective Connie Brueckner served simultaneously as YWCA Board member and lead investigator of complaints against YWCA employees---an unprecedented conflict enabling systematic constitutional violations
\item \textbf{Multiple Victim Accounts:} Documented pattern of client abuse, HIPAA violations, ADA discrimination, and fraudulent court filings affecting numerous families
\item \textbf{Federal Grant Fraud:} Misrepresentation to federal/state agencies while operating with untrained staff, overcrowded facilities, and systematic confidentiality breaches
\item \textbf{DOJ Investigation Context:} Pattern occurs against backdrop of 2012--2016 federal investigation finding Missoula was ``rape capital of the country'' due to institutional failures
\item \textbf{Coordinated Cover-ups:} Strip search allegations at affiliated LifeGuard Group facility (Sheriff's Office Chaplain-run) with no criminal investigation despite DOJ agent concerns
\item \textbf{Estimated Institutional Liability:} \$42.8 million+ across multiple victims (conservative estimate)
\end{itemize}

\subsection*{Law Enforcement Misconduct}

\textbf{Missoula Police Department:}
\begin{itemize}
\item Officer Ethan Smith: Close personal relationship with E'Lise Chard, harassment of Nuno family, removed for ``obvious impropriety''
\item Detective Connie Brueckner: YWCA Board member conflict, warrantless searches, false imprisonment, career destruction tactics
\item Supervisory failures: No meaningful oversight or accountability
\item Pattern and practice: Civil rights violations as institutional policy
\end{itemize}

\textbf{Seattle/Edmonds/King County Agencies:}
\begin{itemize}
\item Interstate coordination with Montana despite no legitimate basis
\item Excessive bail demands (\$50,000--\$100,000) for misdemeanors
\item Evidence suppression and Brady violations
\item OPA complaint falsification (Case 2016OPA-1167)
\end{itemize}

\subsection*{Prosecutorial Misconduct}

\begin{itemize}
\item Charges filed without probable cause across multiple jurisdictions
\item Suppression of exculpatory evidence
\item Witness intimidation and retaliation
\item Malicious prosecution spanning 2015--2020
\item Coordination between Washington and Montana prosecutors
\end{itemize}

\section*{FEDERAL JURISDICTION}

\subsection*{RICO Violations (18 U.S.C. § 1962)}

\textbf{Enterprise:} YWCA Missoula, Missoula Police Department, Missoula County Prosecutor's Office, associated officials

\textbf{Pattern of Racketeering Activity:} 20+ predicate acts spanning 2012--2025, including:
\begin{itemize}
\item Mail/wire fraud (federal grant applications, sponsor reporting)
\item Obstruction of justice (evidence suppression, complaint suppression)
\item Interstate criminal activity (coordination between WA and MT)
\item Witness tampering and intimidation
\item HIPAA criminal violations (42 U.S.C. § 1320d-6)
\end{itemize}

\textbf{Interstate Commerce:} Federal funding, multi-state coordination, wire communications

\textbf{Damages:} \$60+ million across all victims (Nuno: \$18M+; Institutional: \$42.8M+)

\subsection*{Civil Rights Violations}

\begin{itemize}
\item 42 U.S.C. § 1983: Deprivation of rights under color of law
\item 42 U.S.C. § 1985: Civil rights conspiracy
\item First Amendment: Retaliation for protected speech (YWCA complaint, public criticism)
\item Fourth Amendment: Unreasonable searches and seizures, warrantless home invasions
\item Eighth Amendment: Excessive bail and cruel punishment
\item Fourteenth Amendment: Due process violations, malicious prosecution
\end{itemize}

\section*{STATE JURISDICTION}

\subsection*{Montana Criminal Statutes}

\begin{itemize}
\item MCA 45-8-212/213: Criminal defamation
\item MCA 45-7-202: Perjury and false swearing
\item MCA 45-4-102: Criminal conspiracy
\item MCA 45-5-203: Intimidation
\item MCA 45-7-206: Obstructing justice
\item Nonprofit governance violations
\item Professional ethics violations (law enforcement, social work)
\end{itemize}

\subsection*{Washington Criminal Statutes}

\begin{itemize}
\item RCW 9A.72.020: Perjury in the first degree
\item RCW 9A.76.175: Making false statement
\item RCW 9A.72.085: Witness tampering
\item RCW 9A.36.070: Coercion
\item Interstate coordination in criminal enterprise
\end{itemize}

\section*{SCOPE BEYOND NUNO CASE: ENDEMIC INSTITUTIONAL CORRUPTION}

While Elvis Nuno's case provides comprehensive documentation of systematic abuse, \textbf{the evidence raises serious concerns extending beyond reasonable doubt} that warrant expanded criminal investigation into institutional patterns:

\subsection*{Multiple Victim Documentation}

\textbf{Arthur Brown} (Google Review, 2018): YWCA filed defamatory court letters recommending permanent restraining orders without adequate knowledge, separating father from children through fraud upon the court.

\textbf{Anonymous Client \#1} (Facebook, 2020s): HIPAA violations where counseling information was leaked to abusive ex-partner through YWCA staff personal connections; systematic ADA discrimination against disabled clients.

\textbf{Anonymous Client \#2} (Facebook, 2020s): Discriminatory service provision based on client compliance rather than need; fraudulent reporting to federal sponsors claiming to help families while providing substandard care; pattern of repeat clients not properly served.

\textbf{Multiple Employees} (2023--2024 reviews): Documented pattern of ``racist undertones,'' ``untrained staff,'' ``poor management,'' ``exploitative'' compensation, organizational dysfunction creating dangerous conditions for vulnerable clients.

\subsection*{Federal Oversight Failures}

\textbf{2012--2016 DOJ Investigation:} Missoula identified as ``rape capital'' due to systematic sexual assault response failures. YWCA designated as ``community partner'' in reform --- creating additional conflict when YWCA board members investigate YWCA complaints.

\textbf{2023 LifeGuard Group Investigation:} Montana DOJ Agent Andrew Yedinak documented ``concerning'' strip search allegations involving survivor and child at facility run by Sheriff's Office Chaplain Lowell Hochhalter (received \$500,000+ from Gianforte Foundation). No criminal investigation launched; safe houses remain unregulated.

\subsection*{Pattern Recognition}

\begin{itemize}
\item Systematic confidentiality violations endangering domestic violence survivors
\item Federal grant fraud spanning over a decade (\$500,000+ documented)
\item Discrimination against protected classes (disabled, minorities)
\item Coordination with law enforcement to suppress complaints
\item Institutional immunity preventing meaningful accountability
\item Chilling effect preventing victim reporting
\item Ongoing public safety risks to vulnerable populations
\end{itemize}

\section*{REQUESTED ACTIONS}

\subsection*{Federal Investigation}

\begin{enumerate}
\item \textbf{FBI Civil Rights Division:} Comprehensive investigation of institutional conspiracy, pattern and practice violations, interstate criminal activity
\item \textbf{Federal Grand Jury:} Impaneling for RICO prosecution
\item \textbf{HHS Office for Civil Rights:} HIPAA compliance investigation and criminal penalties
\item \textbf{Federal Grant Audit:} DOJ/HHS investigation of grant fraud and misrepresentation
\item \textbf{Asset Forfeiture:} 18 U.S.C. § 1963 forfeiture of RICO enterprise proceeds
\end{enumerate}

\subsection*{State Investigation}

\begin{enumerate}
\item \textbf{Montana Attorney General:} Criminal prosecution of individual perpetrators (Chard sisters) and institutional actors (YWCA officials, MPD officers)
\item \textbf{Washington Attorney General:} Prosecutorial misconduct investigation, interstate coordination analysis
\item \textbf{State Inspectors General:} Comprehensive institutional corruption investigation
\item \textbf{Professional Licensing Boards:} Disciplinary proceedings (law enforcement, social work, legal)
\end{enumerate}

\subsection*{Institutional Reform}

\begin{enumerate}
\item Remove law enforcement officers from YWCA governance
\item Establish independent oversight and ethics committee
\item Implement transparent complaint procedures
\item Mandate federal monitoring (5-year consent decree)
\item Create victim compensation fund
\item Revoke 501(c)(3) nonprofit status pending compliance
\end{enumerate}

\section*{EVIDENTIARY STANDARD}

The submitted documentation establishes criminal conduct \textbf{beyond reasonable doubt} through:

\begin{itemize}
\item Sworn declarations under penalty of perjury
\item Official police reports and court records
\item Documented timeline with precise dates and actors
\item Pattern evidence across multiple victims and timeframes
\item Financial records documenting damages
\item Expert analysis of legal violations
\item Independent corroboration from multiple sources
\item Institutional records (ProPublica nonprofit filings, DOJ reports)
\end{itemize}

\section*{VICTIM IMPACT}

\subsection*{Elvis Ryland Nuno}

\begin{itemize}
\item 11+ years of systematic harassment and false prosecution (2015--present)
\item Complete career destruction in telecommunications field
\item \$2.1 million in lost income and professional advancement
\item \$150,000+ immediate contract losses from warrantless arrests
\item Forced relocation across multiple states fleeing persecution
\item Severe psychological trauma and emotional distress
\item Loss of family relationships and social standing
\item Ongoing threats preventing return to Montana
\item \$6.4--8.4 million in time-barred civil rights claims due to attorney malpractice
\end{itemize}

\subsection*{Broader Community Impact}

\begin{itemize}
\item Vulnerable populations (domestic violence survivors, disabled clients) systematically abused
\item Children separated from parents through fraudulent court filings
\item Federal funding diverted from legitimate services
\item Public trust in law enforcement and victim advocacy destroyed
\item Chilling effect preventing legitimate complaints
\item Pattern of abuse enabling continued criminal activity
\end{itemize}

\section*{URGENCY}

Immediate investigation is critical because:

\begin{enumerate}
\item \textbf{Ongoing Criminal Activity:} Institutional corruption continues to victimize vulnerable populations daily
\item \textbf{Statute of Limitations:} Some predicate acts approaching expiration
\item \textbf{Evidence Preservation:} Risk of destruction as awareness of investigation grows
\item \textbf{Public Safety:} Current YWCA operations at ``full capacity'' with overcrowding and inadequate staffing creating dangerous conditions
\item \textbf{Interstate Coordination:} Multi-jurisdictional criminal activity requires coordinated federal response
\item \textbf{Institutional Immunity:} Local oversight has completely failed; federal intervention essential
\end{enumerate}

\section*{CONCLUSION}

This comprehensive criminal investigation referral documents systematic institutional corruption that extends far beyond the victimization of Elvis Nuno. While his case provides detailed evidence of coordinated criminal enterprise activity by Danielle Chard, E'Lise Chard, and associated institutional actors, the supplemental evidence establishes \textbf{beyond reasonable doubt} that their individual criminal conduct represents the visible manifestation of endemic corruption affecting numerous victims over more than a decade.

\textbf{The pattern includes:}
\begin{itemize}
\item Federal grant fraud and HIPAA violations
\item Systematic civil rights abuses across multiple victim classes
\item Law enforcement misconduct and prosecutorial corruption
\item Institutional capture through governance conflicts
\item Coordinated cover-ups of serious criminal allegations
\item Complete failure of local accountability mechanisms
\end{itemize}

\textbf{Requested Action:} Immediate federal and state criminal investigations of:
\begin{enumerate}
\item Danielle Christine Chard and E'Lise Marie Chard (individual perpetrators)
\item YWCA Missoula and associated governance (institutional enterprise)
\item Missoula Police Department officers and supervisors (law enforcement actors)
\item Missoula County and King County prosecutors (prosecutorial actors)
\item Broader pattern of endemic corruption affecting multiple victims (expanded investigation)
\end{enumerate}

The evidence is comprehensive, well-documented, and meets the threshold for criminal prosecution. Federal RICO charges, civil rights violations, and coordinated state prosecutions are warranted to achieve justice for victims and prevent continued abuse of vulnerable populations.

\vspace{0.3in}

\noindent
Respectfully submitted,

\vspace{0.3in}

\noindent
MISJustice Alliance \\
On behalf of Elvis Ryland Nuno and other victims

\vspace{0.2in}

\noindent
December 19, 2025

\vspace{0.3in}

\noindent
\textbf{Enclosures (11 Sets of Documents \& Relevant Evidentiary Documents):}
\begin{enumerate}
\item Executive Summary / Case Summary.
\item FBI Seattle Field Office Civil Rights Cover Letter (Washington-focused).
\item FBI Montana Field Office Civil Rights Cover Letter (Montana-focused).
\item Montana Attorney General Cover Letter.
\item This Washington Attorney General Cover Letter.
\item U.S. Attorney, District of Montana Cover Letter.
\item U.S. Attorney, Western District of Washington Cover Letter.
\item DOJ Public Integrity Section Cover Letter.
\item Danielle Chard Criminal Report.
\item E'Lise Chard Criminal Report.
\item YWCA Institutional Corruption Supplemental, Elvis Nuno Sworn Declaration, and comprehensive evidentiary documentation.
\end{enumerate}

\end{document}
